\hypertarget{abstract}{%
\subsection{Abstract}\label{abstract}}

Uniform neural reranking may waste computation in retrieval-augmented
generation. We evaluate an interpretable router selecting dense, hybrid,
or hybrid-plus-reranking retrieval for Thai question answering. A frozen
experiment uses 60 held-out questions from 55 article groups, a shared
corpus of 3,223 chunks, and three retrieval timing repetitions. The
router achieves Recall@5 of 0.9833, matching always-reranking, while
reducing mean retrieval latency from 1.0511 to 0.2013 seconds, including
routing overhead. However, nDCG@5 decreases from 0.9314 to 0.9231, and
dense and hybrid baselines remain faster. Removing reranking improves
the router's observed nDCG to 0.9399. Direct mean request latency is
42.489 seconds for the router and 42.141 seconds for always-reranking,
with uncertainty spanning zero difference. Independent
semantic-correctness and faithfulness judgments remain pending. These
results support a limited retrieval-efficiency tradeoff, not general
superiority, answer non-inferiority, or hallucination reduction.

\textbf{Keywords:} Thai question answering; retrieval-augmented
generation; adaptive retrieval; hybrid retrieval; computational
efficiency.

\hypertarget{introduction}{%
\subsection{Introduction}\label{introduction}}

Retrieval-augmented generation combines parametric language models with
external evidence \citep{ref1}. Retrieval can fail through lexical
mismatch, semantic ambiguity, or poorly ranked evidence. Dense passage
retrieval provides a learned semantic representation \citep{ref2}, while
lexical retrieval preserves exact-term signals \citep{ref3}. Combining
these approaches and adding a cross encoder can improve evidence ranking
but introduces additional computation.

A fixed expensive pipeline may be unnecessary for some questions, but an
adaptive pipeline also incurs decision overhead. The relevant research
question is therefore whether query-adaptive selection can reduce
measured costs while preserving answer quality and evidence-grounded
reliability in Thai. This study focuses on an experimentally testable
tradeoff. We do not introduce adaptive retrieval as a new concept.

\hypertarget{research-questions-and-scope}{%
\subsubsection{Research questions and
scope}\label{research-questions-and-scope}}

We distinguish three questions. First, does routing reduce warm
retrieval effort relative to the most expensive fixed pipeline while
preserving retrieval quality? Second, does any retrieval-stage saving
survive once a shared generator produces the answer? Third, do retrieved
passages support the generated claims, independently of whether a short
reference string appears in the response? These questions require
different observations. A faster retrieval call does not establish a
faster complete request, and a relevant passage does not establish a
faithful answer.

The working hypothesis is a conditional tradeoff: selective reranking
can avoid unnecessary cross-encoder work, but may lose evidence quality
when the decision features are weak. A stronger hypothesis that routing
simultaneously improves ranking, answer correctness, faithfulness and
cost is not assumed. We report unfavorable comparisons and ablations
because they identify the boundary of the proposed decision mechanism.
Reliability remains an evaluation objective, rather than an achieved
property asserted by the project title.

The study is deliberately limited to a transparent three-action policy
on one Thai corpus and a modest local model configuration. Its empirical
contribution consists of matched baselines, explicit decision overhead,
a reproducible split and query manifest, and a diagnosis of
disagreements between derived relevance and passage content. This scope
does not establish a first adaptive RAG method or a general optimum for
Thai retrieval. It also does not replace specialist domain evaluation in
law, medicine or other high-stakes applications.

\hypertarget{related-work}{%
\subsection{Related Work}\label{related-work}}

RRF combines rankings without requiring score calibration \citep{ref4},
although subsequent analysis shows that fusion behavior depends on
parameters and tuned score combinations can be stronger \citep{ref5}.
Cross-encoder reranking jointly processes query and passage
\citep{ref6}. These methods motivate the fixed retrieval baselines.

Adaptive-RAG learns to choose retrieval effort from question complexity
\citep{ref7}. Self-RAG integrates retrieval and reflection decisions
into generation \citep{ref8}. MBA-RAG makes cost-sensitive retrieval
choices through a bandit formulation \citep{ref9}. External-feature
adaptive retrieval also predates this study \citep{ref10}. Thus, the
candidate contribution is a controlled Thai evaluation of transparent
A/B/C routing and its overhead, rather than an unrestricted novelty
claim. MIRACL provides a complementary multilingual retrieval benchmark
\citep{ref11}. Ragas motivates separate evaluation of retrieval
relevance and grounded generation \citep{ref12}, but automatic judgments
require Thai-specific validation.

Thai-specific work further constrains novelty. Similarity-based Thai QA
has measured quality and runtime \citep{ref13}. Thai legal RAG has
combined sparse/dense retrieval and neural reranking \citep{ref14}. A
recent Thai Royal Gazette study compares retrieval units and methods
with span-level evidence judgments \citep{ref15}. TSWAP includes
query-dependent forced retrieval for Thai entity lookups \citep{ref16}.
These precedents mean that neither Thai hybrid retrieval nor Thai
query-dependent retrieval can be claimed as new. Our proposed
contribution remains a narrowly controlled evaluation of discrete
strategy selection and its measured costs.

\hypertarget{position-relative-to-earlier-adaptive-systems}{%
\subsubsection{Position relative to earlier adaptive
systems}\label{position-relative-to-earlier-adaptive-systems}}

The action space is a central distinction. Systems that decide whether
to retrieve at all, whether to perform multiple retrieval rounds, or
whether to reflect on an answer address different decisions from
choosing among three retrieval implementations for a single evidence
budget. Our policy always performs an initial dense search. This imposes
a minimum query-encoding cost even when lexical evidence is useful, but
also makes the dense score margin available without a second search. The
comparison therefore evaluates selective additional work after dense
retrieval, not arbitrary selection among independently priced black-box
systems.

External-feature routing is close prior work and limits claims about
novelty \citep{ref10}. Query length and lexical indicators are
inexpensive signals, but cheap observability does not establish that
they predict the marginal benefit of a cross encoder. The present
evaluation consequently treats the features as hypotheses and tests
their removal. We do not infer novelty from using Thai text alone.
Existing Thai hybrid and reranking studies \citep{ref14}, and Thai
query-dependent retrieval \citep{ref16}, already occupy important parts
of that design space.

Evaluation design is equally consequential. MIRACL provides a
multilingual retrieval resource \citep{ref11}, while the Thai Royal
Gazette study uses evidence spans to compare retrieval units
\citep{ref15}. Our source-span conversion is operationally simple but
does not provide exhaustive pooled relevance. A second benchmark with
independently judged evidence would be required to establish whether
observed policy behavior transfers beyond Wikipedia-derived factoid
questions. Ragas \citep{ref12} provides an automated evaluation
framework, but using its terminology or an unvalidated judge would not
supply independent Thai human evidence for this experiment.

\hypertarget{proposed-method}{%
\subsection{Proposed Method}\label{proposed-method}}

Given query q, compute a normalized dense embedding and cosine scores
against a fixed chunk index. Reuse this search if further processing is
selected. Extract Thai-token count, an exact-token proxy indicating
digits/Latin strings/quotation marks, and the score margin between the
first two dense candidates.

A provisional deterministic policy selects reranking when the query is
long or the margin is small; otherwise it selects hybrid retrieval for
exact-token queries and dense retrieval for remaining queries. The fixed
thresholds are 24 Thai tokens and a cosine-score margin of 0.02. They
were frozen before test access and were not calibrated on test outcomes.
Neither token length nor margin is a validated measure of reasoning
complexity or correctness. The selected pipeline returns the same number
of evidence chunks as every baseline.

Hybrid retrieval uses RRF over truncated BM25 and dense rankings.
Reranking applies a shared multilingual cross encoder to the fused
candidate list. The no-hybrid ablation uses dense candidates for
reranking. Feature ablations suppress one decision input at a time. All
query-side work is included in online retrieval timing.

\hypertarget{retrieval-and-fusion-procedure}{%
\subsubsection{Retrieval and fusion
procedure}\label{retrieval-and-fusion-procedure}}

Let D be the fixed passage collection and q the query. The dense encoder
maps the prefixed query and passages to normalized vectors. Passage
scores are dot products, equivalent to cosine similarities under
normalization. Dense retrieval returns the highest-scoring candidates
with a stable tie rule. BM25 operates on lowercased Thai newmm tokens.
Sparse and dense candidate lists each contain at most 20 passages.
Reciprocal-rank fusion assigns each candidate the sum of 1/(60 + rank)
over lists in which it appears. The fused list is truncated to 20
candidates before optional reranking.

The cross encoder scores query--passage pairs jointly and orders the
same candidate set; it does not retrieve unseen passages. Every strategy
returns five passages to generation. Consequently, differences in
answers can arise from which passages appear and their order, although
the passage-count budget is fixed. Character lengths and generator token
counts can still differ across strategies, so equal K does not imply
equal input-token cost. Token counts are measured rather than inferred
from the evidence count.

\hypertarget{transparent-decision-rule}{%
\subsubsection{Transparent decision
rule}\label{transparent-decision-rule}}

Define L as the Thai-token count, E as the presence of a digit, Latin
character or quotation mark, and M as the difference between the top two
dense scores. The policy chooses C when L is at least 24 or M is below
0.02; otherwise it chooses B when E is true, and A in the remaining
cases. The order is consequential: a low-margin exact-token query
selects C before the lexical indicator is considered. The rule does not
use answer labels, source titles, reference strings or generated text as
inputs.

The name ``complexity'' refers only to a length proxy. The
implementation does not estimate reasoning steps, named entities or
semantic ambiguity directly. In particular, a Thai name need not trigger
the Latin/digit/quote indicator, and two overlapping answer-bearing
passages can produce a small dense margin without indicating difficulty.
These limitations are part of the method rather than hidden assumptions
about feature validity.

\hypertarget{ablation-controls-and-implementation}{%
\subsubsection{Ablation controls and
implementation}\label{ablation-controls-and-implementation}}

Each ablation removes exactly one component from the fixed policy.
Removing length, exact-token indicators or margin suppresses the
corresponding condition. Removing hybrid retrieval replaces B by A and
uses dense candidates for C reranking. Removing reranking maps C to B.
These are policy interventions; they need not preserve route frequencies
or candidate sets. We therefore report both quality and actual effort,
and do not treat every ablation as an equal-cost comparison.

The code uses a single shared retrieval implementation for the original
runner and the integrated latency benchmark. Extraction of the retrieval
function changes its callable scope without changing scoring, fusion,
model revisions, thresholds or tie handling. The integrated run records
parity with prior rankings and answers. Random seeds, source hashes,
dataset hashes, package versions and the evaluation schedule accompany
raw outputs. Assertions reject incomplete paired analyses instead of
silently reporting a subset as the planned experiment.

\hypertarget{experimental-setup}{%
\subsection{Experimental Setup}\label{experimental-setup}}

The initial corpus derives from the training partition of the public
iApp Thai Wikipedia QA repository \citep{ref17}. We validate answer
offsets, deduplicate exact normalized questions, and partition connected
components of identical normalized titles or contexts by a deterministic
hash. The fixed retrieval corpus includes all source contexts, while
evaluation labels remain partitioned. Identical contexts cannot cross
partitions. A corpus-only audit identified duplicate-title/context
pairs, which were joined before the frozen evaluation. Three
distinct-entity template pairs remain a near-duplicate sensitivity
concern; model-pretraining contamination remains unknown.

B1, B2, B3 and P1 share multilingual E5 small embeddings, the corpus,
chunking and retrieval budget. B3/P1 share a multilingual mMiniLM cross
encoder. Local generation is implemented using a fixed instruction
model, common Thai prompt and deterministic decoding. Model identities
and resolved revisions are recorded by the runner. The initial generator
is resource-driven and requires Thai suitability validation.

Recall@5, MRR@5 and binary nDCG@5 quantify derived span relevance.
Strict normalized exact match and human semantic correctness evaluate
answers. Independent Thai annotators will label atomic claims for
support and contradiction; abstention evaluation requires validated
unanswerable examples. Wall latency, process CPU time, reranked pairs
and generator token counts quantify efficiency. Bootstrap intervals
resample source-context groups while retaining paired system outputs. A
proposed 0.03 correctness non-inferiority margin was not finalized as a
confirmatory test. No answer non-inferiority claim is made. The present
analysis is exploratory; further tuning must use development/validation
data and a new untouched evaluation sample.

Hardware detected: Intel Core i5-12450H, approximately 32 GiB system
RAM, NVIDIA RTX 3050 laptop-class GPU with 4 GiB VRAM. The initial
retrieval configuration uses CPU inference with four torch threads. This
hardware statement is environment evidence, not a performance
measurement.

\hypertarget{dataset-provenance-and-leakage-controls}{%
\subsubsection{Dataset provenance and leakage
controls}\label{dataset-provenance-and-leakage-controls}}

Only the pinned upstream training file is used to construct the custom
research partitions. Each retained question has a validated source
answer offset and a nonempty set of chunks fully covering that source
span. NFC normalization, zero-width-space removal and whitespace
normalization are applied to text. Chunking uses raw character windows
of 600 characters with 100-character overlap, and identifiers combine
the raw-context hash and offset. This prevents distinct raw whitespace
variants from overwriting one another before normalized grouping.

Questions are deduplicated by normalized question text. Contexts sharing
a normalized title or normalized text are joined into connected
components before a deterministic hash allocates groups to development,
validation or test. Article groups do not cross partitions. The shared
corpus includes contexts associated with every partition, as in
open-book retrieval, while test questions and answer labels do not
select thresholds. Access to test documents is therefore permitted by
design; using test answer labels to fit a router would not be.

The test query manifest contains 60 predetermined questions selected by
seeded hash order, rather than by performance or ease. It covers 55
article groups, so treating all questions as independent would ignore
within-article dependence. The final test partition has additional
unused questions, but evaluating only this fixed sample limits precision
and coverage. The sample is not described as a full benchmark
evaluation. Hash-based grouping reduces known leakage paths but cannot
identify every paraphrase, template similarity or occurrence in
pretrained model data.

\hypertarget{models-and-resource-controls}{%
\subsubsection{Models and resource
controls}\label{models-and-resource-controls}}

The dense model is intfloat/multilingual-e5-small at revision
\texttt{614241f622f53c4eeff9890bdc4f31cfecc418b3}. The reranker is
cross-encoder/mmarco-mMiniLMv2-L12-H384-v1 at revision
\texttt{1427fd652930e4ba29e8149678df786c240d8825}. The generator is
Qwen/Qwen2.5-0.5B-Instruct at revision
\texttt{7ae557604adf67be50417f59c2c2f167def9a775}. These revisions,
rather than mutable model names alone, define the experiment. Downloaded
E5 and generator weight files pass recorded SHA256 checks. The local
generator is chosen for feasibility, not on evidence that it is the
strongest Thai model.

The generation prompt instructs the model to answer in Thai using only
the supplied evidence, cite passage identifiers and abstain when
evidence is insufficient. It also instructs the model not to follow
instructions embedded in evidence. Decoding is greedy with a maximum of
128 new tokens and a 4,096-token request budget. Exceeding the budget
raises an error rather than introducing method-specific truncation. Long
passage identifiers are retained across all systems for consistency;
their cost is a limitation to address only in a future development
protocol.

\hypertarget{measurement-boundaries}{%
\subsubsection{Measurement boundaries}\label{measurement-boundaries}}

The initial retrieval experiment measures query encoding, score
computation, optional BM25/fusion, routing and reranking within each
retrieval call. System order is shuffled for each query and timing
repetition. Corpus encoding and index construction are setup operations,
reported separately. All experiments use local inference, so external
API charges are zero; this excludes electricity and hardware
amortization. CPU process time counts processor time across threads and
is not identical to elapsed wall time.

A follow-up integrated benchmark executes a complete request rather than
reusing saved retrieval outputs. Its timer starts before query encoding
and ends after final answer decoding. The interval contains routing,
retrieval, reranking, evidence lookup and assembly, chat-template
processing, tokenization and generation. Model/index loading and file
logging are outside the request interval. Each system is warmed on a
development query. The frozen test query order is randomized, and
four-query blocks balance each system's position within a query. This
benchmark makes one request per method per question; it cannot estimate
repeated-request jitter for a single question or production queueing
delay.

\hypertarget{statistical-estimands}{%
\subsubsection{Statistical estimands}\label{statistical-estimands}}

Recall@5 is the fraction of labeled relevant chunks retrieved within the
first five positions. MRR@5 uses the reciprocal position of the first
labeled relevant chunk, or zero if absent. Binary nDCG@5 discounts each
labeled hit by its rank and normalizes by the ideal ranking. Unjudged
passages are treated as nonrelevant in these computations; the resulting
incompleteness is investigated rather than concealed. Retrieval timing
repeats are averaged within query before paired comparisons, avoiding
threefold inflation of the quality sample size.

For uncertainty, 2,000 bootstrap replicates sample the 55 article groups
with replacement and retain all associated query pairs. The same
resampled indices are applied to every system. We report percentile
intervals for mean differences in Recall@5, nDCG@5 and latency. Mean,
median and p95 latency are also reported with cluster intervals;
differences of medians and p95 values are distinct from quantiles of
per-query differences. Tail estimates are unstable with only 60
observations. The analysis is exploratory, without multiplicity-adjusted
hypothesis tests, and does not declare statistical significance or
non-inferiority from overlapping or nonoverlapping intervals alone.

\hypertarget{independent-answer-evaluation-protocol}{%
\subsubsection{Independent answer evaluation
protocol}\label{independent-answer-evaluation-protocol}}

A fixed blinded packet contains all 240 original responses, their
questions, reference answers and retrieved evidence, identified only by
opaque response IDs. Method names, routes, timings and retrieval scores
are omitted. A separate coordinator mapping permits linkage after
annotation. Two Thai-proficient raters should independently judge
semantic correctness and segment factual claims, recording supporting
passage IDs and rationales. Independent raters are not yet available;
all human labels therefore remain missing, not correct, incorrect or
faithful by default.

Correctness permits paraphrase and equivalent values but rejects a
materially wrong answer or misleading additional assertions. Ambiguous
questions and defective references receive an uncertainty flag for
adjudication. Claim support distinguishes entailed, contradicted,
unsupported and uncertain claims. A factually correct claim can be
ungrounded in the supplied passages, and a claim faithfully copied from
the wrong entity's passage can fail to answer the intended question.
Citations alone do not establish support. Unsupported claims combine
contradicted and unsubstantiated claims for the aggregate while
retaining the subtypes in raw ratings.

Before adjudication, report agreement and Cohen's kappa where its
denominator is defined. Preserve disagreements and claim-boundary
decisions. Faithfulness is supported factual claims divided by all
adjudicated factual claims, with annotation coverage stated. Zero-claim
refusals have undefined claim faithfulness and do not count as perfect
responses. Response-level hallucination is operationalized as at least
one unsupported factual claim among adjudicated factual responses. The
answerable-only dataset cannot support balanced abstention accuracy.
These definitions are a protocol, not collected measurements.

\hypertarget{results}{%
\subsection{Results}\label{results}}

The fixed test sample contains 60 questions from 55 article groups. All
nine systems (four primary systems and five ablations) completed three
repetitions, producing 1,620 raw retrieval records. Query IDs match the
frozen list; group disjointness and identical rankings across
repetitions were checked. Index preparation required 158.07 seconds and
is excluded from warm retrieval timing.

\begin{longtable}[]{@{}lrrrr@{}}
\toprule
System & Recall@5 & MRR@5 & nDCG@5 & Mean retrieval (ms)\tabularnewline
\midrule
\endhead
B1 dense & 0.9500 & 0.9139 & 0.9232 & 17.49\tabularnewline
B2 hybrid & 0.9667 & 0.9131 & 0.9263 & 25.10\tabularnewline
B3 hybrid + reranking & 0.9833 & 0.9139 & 0.9314 &
1051.13\tabularnewline
P1 adaptive & 0.9833 & 0.9028 & 0.9231 & 201.29\tabularnewline
\bottomrule
\end{longtable}

P1 reduces observed mean retrieval time by 80.85\% relative to B3 and
reranks 3.33 rather than 20 pairs per query (83.33\% fewer). Mean router
overhead is 0.230 ms and is included. P1 routes 48 questions to dense,
two to hybrid, and ten to reranking. Its nDCG difference against B3 is
−0.00833 (article-cluster bootstrap 95\% interval {[}−0.02343,
0.00000{]}); the retrieval-time difference is −0.84984 seconds
{[}−0.94705, −0.74234{]}. P1 is approximately 11.51 times slower than B1
and 8.02 times slower than B2. Its observed nDCG is lower than both,
although Recall@5 is higher. Thus no primary method is universally best
across the selected objectives.

\begin{longtable}[]{@{}lrrr@{}}
\toprule
Removed P1 component & Recall@5 & nDCG@5 & Mean retrieval
(ms)\tabularnewline
\midrule
\endhead
Query length & 0.9833 & 0.9231 & 174.39\tabularnewline
Exact-token indicator & 0.9833 & 0.9231 & 193.41\tabularnewline
Dense score margin & 0.9500 & 0.9232 & 52.02\tabularnewline
Hybrid retrieval & 0.9500 & 0.9068 & 195.63\tabularnewline
Reranking & 0.9667 & 0.9399 & 19.25\tabularnewline
\bottomrule
\end{longtable}

Removing reranking yields the highest observed nDCG among the tested
configurations, while reducing recall relative to full P1. This is an
exploratory ablation result, not a newly selected confirmatory policy.
Length and exact-token removals do not change aggregate retrieval
quality on this sample. The score-margin feature improves recall but
does not establish an nDCG benefit.

Data preparation retained 5,667 questions (3,360 development, 1,133
validation, 1,174 test) and excluded 94 records for invalid/uncovered
spans or duplicate questions. All corpus chunks are available as an
open-book index; test query labels do not train or tune the router. The
60-question test sample is a subset, not the entire held-out partition.

All 240 matched generation responses completed. Strict exact match is
0.0333 for B1, B2 and P1, and 0.0167 for B3. These low scores are not
semantic correctness estimates because the metric penalizes explanatory
text and citations. Estimated mean request latency is 42.275 seconds for
B1, 41.758 for B2, 42.841 for B3 and 43.495 for P1. P1 therefore
provides no observed overall latency saving; its paired difference from
B3 is +0.6545 seconds {[}−0.5456, 1.8908{]}. P1 uses 1,940.22 input and
60.82 output tokens per query, compared with B3's 1,917.45 and 54.58.
Total generation usage is 461,175 input and 13,557 output tokens. API
charges are zero because inference is local; this does not imply zero
computational cost. Independent semantic correctness, faithfulness,
unsupported-claim rate, hallucination rate and balanced abstention
accuracy remain pending. Separately measured retrieval and generation
durations provide an estimated request latency; direct full-request
timing is reported below; electricity/hardware monetary cost remains
unmeasured.

\hypertarget{direct-end-to-end-latency-validation}{%
\subsubsection{Direct end-to-end latency
validation}\label{direct-end-to-end-latency-validation}}

The integrated benchmark completed 240 requests on the same 60 frozen
questions, with one request per method/question and balanced
within-query execution positions. The reported interval starts before
dense query encoding and ends after final answer decoding. No separate
stage durations are added to obtain this measurement.

\begin{longtable}[]{@{}lrrrr@{}}
\toprule
Method & Mean (s) & Median (s) & p95 (s) & SD (s)\tabularnewline
\midrule
\endhead
B1 & 41.765 & 41.839 & 53.509 & 7.312\tabularnewline
B2 & 40.590 & 42.051 & 54.061 & 8.439\tabularnewline
B3 & 42.141 & 43.027 & 52.882 & 7.370\tabularnewline
P1 & 42.489 & 41.660 & 55.272 & 7.823\tabularnewline
\bottomrule
\end{longtable}

Mean latency 95\% article-cluster bootstrap intervals are B1 {[}39.911,
43.494{]} seconds; B2 {[}38.481, 42.615{]} seconds; B3 {[}40.272,
43.974{]} seconds; P1 {[}40.545, 44.450{]} seconds.

P1 versus B3 has an observed direct mean latency reduction of -0.83\%
(negative means slower). This comparison measures the complete warm
request and should not be confused with the 80.85\% retrieval-only
reduction.

The P1-minus-B3 mean difference is 0.3477 seconds, with 95\% interval
{[}-1.2076, 1.8974{]}. No formal significance or equivalence test is
claimed.

Rankings match the original run on 240/240 requests, and decoded answers
match on 240/240 requests. Model revisions, thresholds, CPU thread
count, prompt and generation budget remain fixed. The generator's
recorded dtype is bfloat16 on CPU. The workstation was not isolated from
other activity; the benchmark quantifies this session, not guaranteed
production latency.

Median and p95 intervals and all six pairwise comparisons are retained
in the supplementary statistical output. The p95 is sensitive to a few
observations and only one request was timed per query/method.

\hypertarget{discussion}{%
\subsection{Discussion}\label{discussion}}

The observed policy saves work relative to always-reranking but is not
competitive with the cheapest baselines on nDCG and latency jointly. Six
queries incur positive regret under source-derived relevance labels:
four choose reranking and two remain on dense retrieval. These are
label-defined discrepancies, not six independently established semantic
routing errors. Evidence inspection shows that the park query retrieves
an unjudged first passage that explicitly contains the answer. The
band-nationality query also omits the band name. Low score margin does
not reliably identify improvement in the measured ranking metric, but
incomplete judgments and question ambiguity constrain that
interpretation. The no-reranking ablation reinforces this conclusion.
Because relevance is non-exhaustive, these ranking failures should also
be inspected against full passage text.

Qualitative inspection identifies generation errors despite available
evidence: one answer gives the English name of Aphrodite when the
question asks for her Roman name, although the retrieved passage
explicitly gives Venus. Another supplies an unsupported death year.
These assistant-reviewed cases are not independent human annotations and
do not estimate hallucination rates. Long evidence identifiers can
consume the fixed generation budget, and strict exact match penalizes
otherwise relevant explanatory answers. Neither prompt nor router is
retuned after examining these test outcomes.

The defensible contribution is a reproducible, overhead-inclusive Thai
comparison showing where a transparent three-way policy does and does
not offer value. Learned routing is not yet justified: this small sample
provides error hypotheses, not evidence of generalizable predictability.
Calibration or learning should proceed only on development/validation
data, with any new claims evaluated on untouched questions. \#\#\#
Mechanism-level error analysis

Among ten C-route questions, removing reranking improves nDCG on four,
reduces it on one and leaves it unchanged on five. The other 50
questions retain their original action and ranking. The four gains are
0.369070 for the regional fish name, 0.130930 for the abbreviated
Tolkien death question, 0.369070 for park type and 0.569323 for band
nationality. One director question loses 0.430677 because reranking had
recovered the labeled passage at rank four while hybrid missed the top
five. Dividing their net change by 60 yields the observed increase of
0.016795. The one lost hit explains the recall decrease of 1/60.

The fish example is an answer-localization failure: the reranker
promotes the article's general species description over its later
regional-name sentence. The Tolkien example combines an abbreviated
entity with competing death-related passages, so semantic topical
similarity does not ensure the correct subject. The docusate query is a
missed reranking opportunity under the labels: the dense margin is
relatively large, but the answer-bearing secondary use is ranked below
the general introduction. The Diana concert query promotes a similar
sentence about another performance within the same article; the
exact-token proxy does not recognize a Thai personal name as an entity
feature.

The park example changes how the aggregate must be interpreted. Its
first reranked passage explicitly contains the requested park
classification, but the qrel identifies a different source chunk. Thus
an apparent nDCG loss can coexist with adequate evidence. The band
question lacks the band's name, and the director question similarly
omits the work being directed. These source-context-dependent questions
are weak standalone retrieval tests. We preserve their scores and
disclose the problem rather than remove unfavorable examples after
seeing results.

\hypertarget{retrieval-savings-and-complete-request-cost}{%
\subsubsection{Retrieval savings and complete-request
cost}\label{retrieval-savings-and-complete-request-cost}}

A reduction in cross-encoder calls can be large relative to retrieval
time while small relative to generation. Furthermore, different evidence
changes prompt length and output length even when the number of passages
is fixed. The original generation run uses more input and output tokens
for P1 than for B3 on average. Adding separate measurements suggested no
overall saving, but that estimate is not a substitute for a complete
timed request. The integrated benchmark is designed to distinguish these
quantities under the same policy, model and evidence budget.

The system-level objective should therefore be explicit. If the priority
is recovering a source-labeled passage while avoiding always-reranking,
the observed P1 recall/time tradeoff may be useful. If the priority is
maximizing nDCG at minimal latency, dense, hybrid or the no-reranking
ablation can be preferable in this sample. If the objective is reliable
generated answers, neither comparison is decisive without grounded
answer judgments. A contribution that states these boundaries is more
defensible than selecting whichever metric makes the proposed system
look strongest.

\hypertarget{development-only-next-steps}{%
\subsubsection{Development-only next
steps}\label{development-only-next-steps}}

Future work should first improve evaluation validity through pooled
judgments and standalone-question assessment on development data. It may
then estimate the incremental quality benefit of reranking, conditioned
on observable query and retrieval features, under an explicit cost
budget. A calibrated decision tree or other lightweight model would be
justified only by out-of-group development/validation performance and
sufficient labeled decision variation. The six test diagnostics are not
a training set for a new policy. Any new thresholds, entity features,
prompt format or learned router must be frozen before evaluation on a
new untouched sample.

\hypertarget{limitations}{%
\subsection{Limitations}\label{limitations}}

The current evidence labels are derived from answer spans and are not
exhaustive relevance judgments. Overlapping chunks can inflate the
number of relevant items. Character windows may split Thai words.
Wikipedia factoid QA does not represent all Thai information needs.
Model pretraining may include benchmark content. Heuristic features
require calibration and may not transfer. Human grounding evaluation,
residual near-duplicate sensitivity analysis and full literature detail
verification remain incomplete. Three retrieval repetitions quantify
within-session variability, but do not establish cross-machine or
cross-session latency stability. The small held-out sample and 55
article groups limit precision. The small initial generator may be
inadequate for Thai. Generation was run once per query/system; five
responses resumed after an interruption, and session drift may affect
timing. Long evidence IDs and the output cap can also impair answer
formatting. Official conference pages disagree on the submission
deadline and page limits; final formatting requires reconciliation.

\hypertarget{threats-to-validity-and-reproducibility-boundaries}{%
\subsubsection{Threats to validity and reproducibility
boundaries}\label{threats-to-validity-and-reproducibility-boundaries}}

Construct validity is limited by source-span relevance, strict answer
matching and the use of length as a complexity proxy. The park case
shows that derived qrels can penalize a passage that answers the
question. Exhaustive judgments are unavailable, so retrieval scores
measure recovery of labeled source evidence rather than all acceptable
evidence. Exact match can reject correct explanatory responses and
should not be presented as semantic accuracy. Conversely, a matching
reference substring would not prove that all generated claims are
supported.

Internal validity depends on preserving common models, candidate
budgets, prompt instructions and query sets. Fixed revisions and output
hashes make silent changes detectable. Randomized or balanced execution
order reduces some order effects, but this is a workstation experiment
rather than an isolated serving environment. Background load, CPU
frequency, thermal conditions and memory residency can influence
latency. One integrated request per method/query and three retrieval
repetitions do not establish robustness across days, concurrent
workloads or deployment stacks. Confidence intervals quantify the
resampling variability in this sample, not every source of measurement
error.

External validity is restricted to Thai Wikipedia-derived questions, a
small local generator, the selected multilingual retriever/reranker and
CPU hardware. Domain terminology, multi-hop reasoning, long documents
and user-authored queries may behave differently. The derived test
sample contains no independently validated unanswerable subset. The same
public text may have appeared during model pretraining; hash-based
research splits do not resolve that contamination. Results should not be
generalized to all Thai RAG systems or interpreted as evidence that a
particular architecture is unsuitable for other applications.

Conclusion validity is limited by the small sample, dependent article
questions, exploratory ablations and unavailable independent answer
judgments. The highest observed ablation score is not an independently
validated new method, because it was identified after test access. We
report effect sizes and intervals rather than a favorable binary
significance narrative. The correct inference from missing faithfulness
annotations is that reliability remains unestablished. Negative results
and measurement limitations remain part of the record even if future
policies improve.

The reproducibility package stores the source-data revision and hash,
model revisions and weight checks, frozen query IDs, configuration
files, raw rankings and answers, per-request timings, token counts,
environment locks and analysis scripts. Existing held-out artifacts are
hashed before follow-up work and are not overwritten. Blinded annotation
templates contain no judgments, and mapping files are separated from
rater material. Model licenses, corpus attribution and contributor-data
minimization must be respected when sharing the package; source
payment/contributor fields are not used in processed experimental
records.

\hypertarget{conclusion}{%
\subsection{Conclusion}\label{conclusion}}

The frozen policy reduces retrieval latency relative to always-reranking
while matching observed Recall@5, but loses nDCG and remains
substantially slower than dense or hybrid retrieval. Removing reranking
produces a stronger observed nDCG/latency tradeoff, revealing a weakness
in the current decision mechanism. These results support a bounded
empirical contribution and an honest negative finding about the policy,
rather than a claim that adaptive routing is generally superior. Direct
request timing shows no observed mean latency saving for P1 against any
primary baseline; its difference from B3 has a bootstrap interval
spanning zero. Reliability and semantic-answer conclusions await
independent annotation.
